\documentclass{article}

\usepackage{cstyle}

\title{Espsecificaciones para las mejoras de DSR}
\author{
        Carlos Madrid Espinosa \\
        \href{mailto:carlosmadrid@uma.es}{carlosmadrid@uma.es}
}

\begin{document}

\maketitle

\section*{Versiones}\label{sec:ver}

\begin{table}[h]\label{tab:ver-versions}
        \begin{tabular}{ccll}
                \toprule
                Version & Iteración & Lista de Cambios & Fecha  \\
                \midrule
                01      & 00        & Version inicial. & \today \\
                \bottomrule
        \end{tabular}
        \caption{Versiones del documento}
\end{table}

\tableofcontents

\newpage
%\twocolumn

\section{Introducción}\label{sec:intro}

Este documento recoge las necesidades detectadas para la actualización del
sistema DSR (Deep State Representation). Estas se usará como base para el
planteamiento de una base de rerquisitos.

\section{Planteamiento}\label{sec:scope}

El Grupo de Ingeniería de Sistemas Integrados de la Universidad de Málaga
cuenta con diversos sistemas, como robots asistenciales y entornos
inteligentes, que comparten una necesidad común: disponer de una representación
del entorno que pueda utilizarse internamente para la toma de decisiones y la
monitorización. Cada sistema agrupa diferentes elementos de los que es
necesario registrar información. Por ejemplo, en el caso de un robot
asistencial que interactúa con una persona, es preciso registrar datos
específicos de esta (como su identificación), tener en cuenta variables
ambientales (hora, temperatura, humedad, etc.) y registrar la propia
interacción entre la persona y el robot para poder generar acciones adecuadas.

Para dar respuesta a esto, se emplea DSR (Deep State Representation), un
sistema de modelado de información basado en estructurar los componentes del
sistema y sus relaciones en forma de grafo, compuesto por nodos y aristas. De
este modo, cada elemento se representa como un nodo con una serie de atributos
que lo caracterizan, mientras que las relaciones se definen mediante las
aristas, las cuales también disponen de atributos para especificar mejor dicha
conexión.

Un aspecto fundamental de DSR es que actúa puramente como un modelo; por sí
solo, ni genera ni almacena la información por cuenta propia. Para ello,
proporciona una API que cualquier sistema puede utilizar para leer o escribir
en la representación del entorno. Estos componentes interactivos se denominan
agentes, los cuales pueden asumir roles de proveedores de información o de
consumidores de la misma.

Para potenciar la independencia de los agentes, DSR presenta una arquitectura
distribuida en la que cada agente almacena una copia local del estado. Cuando
un agente desea realizar un cambio, envía la diferencia (delta) a un nodo
central, el cual se encarga de propagar la actualización al resto de agentes.
Gracias a este diseño, cada agente lee directamente de su copia local, lo que
reduce drásticamente el número de peticiones necesarias y permite que el
sistema continúe operando con normalidad incluso si se interrumpe temporalmente
la conectividad.

\section{Problemas detectados}\label{sec:problems}

El sistema actual presenta dos problemas fundamentales derivados principalmente
de la implementación \textit{adhoc} actual del modelo de grafo:

\begin{itemize}
        \item \textbf{Velocidad:} La implementación actual del grafo es una
              solución \textit{adhoc} excesivamente compleja que penaliza el
              rendimiento, aportando poca flexibilidad y resultando demasiado
              lenta en la práctica, lo que desincentiva su uso en muchos
              escenarios.
        \item \textbf{Pérdida de datos y divergencia:} Al actualizar al resto de
              sistemas, se producen episodios donde algún dato se pierde.
              Esto provoca que el resto de componentes descarten la
              actualización, haciendo que la información de DSR en los distintos
              agentes derive y se desincronice con el tiempo.
\end{itemize}

\section{Propuesta}\label{sec:proposition}

Para resolver los problemas de rendimiento y complejidad del modelo de grafo,
se propone \textbf{sustituir la implementación \textit{adhoc} actual por una
        librería de gestión de grafos ya existente y contrastada}, mitigando así los
cuellos de botella y ganando flexibilidad.

En cuanto a la pérdida de datos y la sincronización, se plantean tres enfoques
arquitectónicos principales:
\begin{enumerate}
        \item \textbf{Copias distribuidas completas:} Mantener las copias del
              DSR en cada agente, comunicándose todos entre sí de forma
              totalmente mallada.
        \item \textbf{DSR Centralizado (Cliente/Servidor):} Hacer que los
              agentes lean y escriban directamente de un DSR centralizado,
              organizándose en una topología en estrella.
        \item \textbf{Solución mixta:} Combinar agentes con copia local y
              agentes ligeros sin copia. Los agentes con DSR se mantienen
              conectados entre sí, mientras que los agentes sin copia consultan
              a cualquiera de los nodos que sí disponen de ella.
\end{enumerate}

El \textbf{tercer enfoque (solución mixta)} se postula como la mejor opción, ya
que proporciona la máxima flexibilidad operativa adaptándose mejor a las
necesidades reales de los diferentes subsistemas.

Adicionalmente, para mitigar los problemas de divergencia de datos entre
copias, se propone:
\begin{itemize}
        \item \textbf{Mecanismos de integridad (CRC):} Introducir sumas de
              verificación o códigos de redundancia cíclica (CRC) en las
              estructuras del sistema para detectar alteraciones o
              discrepancias de forma temprana.
        \item \textbf{Agente de supervisión de red:} Incorporar un mecanismo
              basado en un agente especial de control cuya función específica
              sea monitorizar la red, validar la consistencia de los grafos y
              detectar posibles derivas de información entre nodos.
\end{itemize}

\section{Requisitos de Portabilidad y Embedabilidad}\label{sec:portability}

Dado que el sistema va a desplegarse en una amplia variedad de plataformas y
dispositivos heterogéneos, se establecen como pilares fundamentales de diseño
los siguientes aspectos:

\begin{itemize}
        \item \textbf{Portabilidad:} El código debe ser altamente portable,
              facilitando su compilación y ejecución en diferentes entornos
              mediante el uso de cadenas de herramientas basadas en LLVM
              (tales como \texttt{clang} o \texttt{clangd}) y el apoyo en
              librerías estándar multisistema compatibles con POSIX.
        \item \textbf{Embedabilidad y eficiencia:} Con el fin de permitir su
              integración en sistemas con recursos limitados o embebidos,
              la arquitectura debe mantenerse lo más ligera y mínima posible,
              priorizando en todo momento la eficiencia computacional y el bajo
              consumo de memoria.
\end{itemize}

\end{document}

